% SGA 3, Expose VI_B, Loop 1 English reconstruction.
% Scope: Corollary 6.2.6 and its proof in full, including source
% editor's note 59.
% Authority: Polo--Gille Exp6B-13oct24.pdf, component-local pages
% 43--44 (combined-reader pages 384--385; re-edition headers
% 373--374; original-margin printed page 359 begins on local 44).
% Remark 6.3 begins on local page 44 / combined page 385 and is
% excluded.
% OCR and Reinhold English Markdown are comparison witnesses only.
% This unit contains no diagrams.

\medskip
\phantomsection\label{cor:sga3-VIB-6.2.6}%
\noindent\textbf{推论 6.2.6.---}%
\footnote{\label{note:sga3-VIB-59}编者注: 我们在此插入
本推论 (参见 \SGARefLink{sga3:xii:6.1}{第~XII 篇, 6.1}), 因为后文将用到它.
在 \SGARefLink{rem:sga3-VIB-6.3}{6.3} 中我们回到第~VI\(_B\) 篇的
原文.}
\emph{设 \(k\) 为域, \(G\) 为连通
代数 \(k\)-群. 则
\(\underline{\operatorname{Centr}}_G(G)\) 可由
\(G\) 的闭子群概形 \(Z\) 表出, 且 \(G/Z\) 是一个仿射
代数 \(k\)-群.}

\emph{证明.}
第一个断言当然已包含于
\hyperref[cor:sga3-VIB-6.2.5]{6.2.5\textup{(i)}} 中, 但我们将
看到它也由第二个断言的证明得到.
事实上, \(G\) 通过伴随表示作用于有限维
\(k\)-向量空间
\[
  V_n=\mathcal O_{G,e}/\mathfrak m_e^{\,n+1},
\]
其中 \(\mathfrak m_e\) 为
\(\mathcal O_{G,e}\) 的极大理想. 设 \(K_n\) 为态射
\[
  \rho_n:G\longrightarrow\operatorname{GL}(V_n).
\]
的核. 由 \hyperref[prop:sga3-VIA-5.4.1]{VI\(_A\), 5.4.1},
\(\rho_n\) 诱导闭浸入
\[
  G/K_n\lhook\joinrel\longrightarrow\operatorname{GL}(V_n),
\]
所以每个 \(G/K_n\) 均仿射. 由于 \(G\) 为 Noether 概形,
\(K_n\) 的交 \(K\) 等于其中某一个 \(K_n\);
因而 \(G/K\) 仿射.

另一方面, 若 \(Z\) 表示 \(G\) 的中心, 则显然
\[
  Z\subset K.
\]
下面证明反向包含关系. 设
\(\widehat{\mathcal O}_{G,e}\) 为
\(\mathcal O_{G,e}\) 关于 \(\mathfrak m_e\)-进拓扑的完备化,
并令
\[
  \widehat S=\operatorname{Spec}\widehat{\mathcal O}_{G,e},
  \qquad
  S=\operatorname{Spec}\mathcal O_{G,e}.
\]
由于 \(\widehat S\to S\) 忠实平坦, 且两个
态射
\[
  K\times_k S\longrightarrow S,
  \qquad
  (g,x)\longmapsto gxg^{-1}
  \quad\text{和}\quad
  (g,x)\longmapsto x
\]
在基变换 \(\widehat S\to S\) 后相同,
故它们相同. 因此 \(K\) 平凡作用于 \(\mathcal O_{G,e}\).

由 \hyperref[ex:sga3-VIB-6.2.4-e]{6.2.4}\textup{(e)}, \(G\) 中被 \(K\)
固定的点所成的子对象 \(G^K\), 即
\(\underline{\operatorname{Centr}}_G(K)\), 是 \(G\) 的闭
子概形. 因此它由 \(\mathcal O_G\) 的一个拟凝聚
理想 \(\mathcal I\) 定义. 由于 \(G^K\) 包含
\[
  S=\operatorname{Spec}\mathcal O_{G,e}
\]
且 \(\mathcal I\) 有限型 (因为 \(G\) 为 Noether
概形), 存在 \(e\) 的一个开邻域 \(U\), 使得
\[
  \mathcal I|_U=0.
\]
因此子群
\[
  G^K=\underline{\operatorname{Centr}}_G(K)
\]
包含 \(U\), 从而也包含 \(U\cdot U\), 而后者等于
\(G\), 因为 \(G\) 不可约
(\hyperref[sga3:VIA:0.5]{VI\(_A\), 0.5}). 因此,
\[
  \underline{\operatorname{Centr}}_G(K)=G.
\]
所以 \(K\subset Z\), 于是 \(Z=K\), 得证.

% 翻译统计: 源行数 98; 输出行数 98 (含本统计行); 行内/行间数学片段 60; 翻译占位 0. 底本疑点另报主线程.
